% GENERATED FILE. Edit canonical lessons, then run npm run generate:books.
% canonical-source-hash: fnv1a64:29b4003a0a63a28e
% canonical-lessons: ZH-C411-matou, ZH-C411-chuanzhang, ZH-W411-chuan-boat, ZH-C411-kaichuan, ZH-C411-duchuan, ZH-C411-jiushengyi

\chapter{On the Water}
\label{ch:zh-on-the-water}
\hlchaptermodality{411}

\begin{chapteropening}
\textbf{By the end of this chapter:} \emph{I can talk about the pier, the captain, a boat setting sail, the ferry and life jackets.}

Complete the last lesson of chapter 411: I can talk about the pier, the captain, a boat setting sail, the ferry and life jackets.
\end{chapteropening}

\section[mǎtou]{mǎtou — pier; wharf}

\label{lesson:ZH-C411-matou}



\begin{quote}
Before the new one: say the Mandarin for to launch a new transport line, then the Mandarin for to hit the road.
\end{quote}

\subsection*{You'll want to know: mǎtou}
\emph{mǎtou} — \textquotedblleft{}pier; wharf\textquotedblright{}.

\begin{quote}
\emph{qù mǎtou zěnme zǒu} — How do I get to the pier?
\end{quote}

Its characters come later; for now it is a word for the ear.

\begin{cousinweb}
\emph{mǎtou} is built from \emph{mǎ} and the noun ending \emph{tou}, both inherited from Middle Chinese, \emph{maeX} and \emph{duw} (Guangyun, in Baxter's transcription). The dictionary lists the first only in the name of agate, a stone like jade, and the second as 'the head'; how the pair came to mean a pier is not settled.
\end{cousinweb}

\begin{grammarlens}[title={What you've built}]
One more word for getting around.
\end{grammarlens}

\subsection*{Your turn}
\begin{itemize}
\raggedright
  \item \emph{Say it:} \emph{mǎtou}
  \item \emph{Say it:} \emph{mǎtou}, once more
  \item \emph{Say it:} \emph{qù mǎtou zěnme zǒu} — How do I get to the pier?
  \item \emph{Recall:} say the Mandarin for to launch a new transport line, then the Mandarin for to hit the road, then say \emph{mǎtou} again
\end{itemize}

\begin{tcolorbox}[breakable,colback=teal!4,colframe=teal!35!black,title={Before you move on}]
What does mǎtou mean? (\textquotedblleft{}Pier; wharf\textquotedblright{}.) Say it once more.
\end{tcolorbox}
\section[\texorpdfstring{\zh{船长} (chuánzhǎng)}{船长 (chuánzhǎng)}]{\zh{船长} (chuánzhǎng) — ship's captain}

\label{lesson:ZH-C411-chuanzhang}



\begin{quote}
Before the new one: say the Mandarin for to hit the road, then the Mandarin for a pier for boats.
\end{quote}

\subsection*{You'll want to know: \zh{船长}}
\textbf{\zh{船长}} — \emph{chuánzhǎng} — \textquotedblleft{}ship's captain\textquotedblright{}.

\begin{quote}
\textbf{\zh{我想当船长}} \emph{wǒ xiǎng dāng chuánzhǎng} — I want to be a ship's captain.
\end{quote}

The new character, in pieces: \zh{船} — eleven strokes in two pieces: on the left the boat piece, six strokes, a short stroke falling to the left, a long one down the left, a bar that turns down and hooks, a dot, a bar rising across and a last dot; on the right a short stroke falling to the left and a bar that turns down and bends right, then \zh{口} under them. The next lesson puts it in your hand.

\begin{cousinweb}
\textbf{\zh{船长}} (\emph{chuánzhǎng}) is built from \zh{船} 'a boat' and \zh{长} 'the head', both inherited from Middle Chinese, \emph{zywen} and \emph{trjangX} (Guangyun, in Baxter's transcription), which glosses the second 'big'.
\end{cousinweb}

\begin{grammarlens}[title={What you've built}]
One more word for getting around.
\end{grammarlens}

\subsection*{Your turn}
\begin{itemize}
\raggedright
  \item \emph{Say it:} \emph{chuánzhǎng}
  \item \emph{Say it:} \emph{chuánzhǎng}, once more
  \item \emph{Say it:} \emph{wǒ xiǎng dāng chuánzhǎng} — I want to be a ship's captain.
  \item \emph{Recall:} say the Mandarin for to hit the road, then the Mandarin for a pier for boats, then say \emph{chuánzhǎng} again
\end{itemize}

\begin{tcolorbox}[breakable,colback=teal!4,colframe=teal!35!black,title={Before you move on}]
What does \zh{船长} mean? (\textquotedblleft{}Ship's captain\textquotedblright{}.) Say it once more.
\end{tcolorbox}
\section[\texorpdfstring{\zh{船} (chuán)}{船 (chuán)}]{\zh{船} — the boat piece, then \zh{口} under a small frame}

\label{lesson:ZH-W411-chuan-boat}



\begin{quote}
The word first, then the character inside it.

\begin{itemize}
\raggedright
  \item \emph{Recall:} say \emph{chuánzhǎng} — \textbf{R1}, one lesson back, before it decays
\end{itemize}

\begin{quote}
\zh{口} \zh{长}
\end{quote}

\emph{Read it:} each one aloud, with nothing said first — \emph{kǒu}, \emph{cháng}
\end{quote}

\subsection*{Script — the boat piece, then \zh{口} under a small frame}
\hlblockfigure{\detokenize{figures/ZH-W411-chuan-boat-filmstrip.pdf}}{How \zh{船} is written, stroke by stroke}

\begin{quote}
\zh{船}
\end{quote}

\textbf{Two pieces, left before right.} Count the pieces, not the strokes.

\textbf{The left piece, the boat piece, six strokes.} A short stroke falling to the left. A long stroke down the left side, falling to the left at the foot. A bar that turns down the right side and hooks to the left. A dot inside, a bar rising to the right across the middle, and a last dot under it.

\textbf{The right piece, five strokes.} A short stroke falling to the left, and a bar that turns down and bends to the right at the foot: a small frame. Then \zh{口} under it, three strokes.

6 + 5 = eleven.

\subsection*{Your turn}
\begin{itemize}
\raggedright
  \item \emph{Write it:} \zh{船}, all eleven, with the model in front of you
  \item \emph{Cover:} write it again from memory, and count the pieces, not the strokes
\end{itemize}

\begin{tcolorbox}[breakable,colback=teal!4,colframe=teal!35!black,title={Before you move on}]
How many pieces does \zh{船} have? (\textbf{Two}.) What is the left piece? (\textbf{The boat piece, six strokes}.) What sits at the foot of the right piece? (\textbf{\zh{口}}.)

Source: \href{https://raw.githubusercontent.com/chanind/hanzi-writer-data/68d10a4b21150cae5e1ebbd223eed289cf32d90c/data/\%E8\%88\%B9.json}{Hanzi Writer Data \zh{船}, pinned stroke paths and medians}.
\end{tcolorbox}
\section[\texorpdfstring{\zh{开船} (kāichuán)}{开船 (kāichuán)}]{\zh{开船} (kāichuán) — to set sail}

\label{lesson:ZH-C411-kaichuan}



\begin{quote}
Before the new one: say the Mandarin for a pier for boats, then the Mandarin for a ship's captain.
\end{quote}

\subsection*{You'll want to know: \zh{开船}}
\textbf{\zh{开船}} — \emph{kāichuán} — \textquotedblleft{}to set sail\textquotedblright{}.

\begin{quote}
\textbf{\zh{船长说十点开船}} \emph{chuánzhǎng shuō shí diǎn kāichuán} — The captain says we sail at ten.
\end{quote}

\begin{cousinweb}
\textbf{\zh{开船}} (\emph{kāichuán}) is built from \zh{开} 'to start off' and \zh{船} 'a boat', both inherited from Middle Chinese, \emph{khoj} and \emph{zywen} (Guangyun, in Baxter's transcription), which glosses the first 'to open'.
\end{cousinweb}

\begin{grammarlens}[title={What you've built}]
One more word for getting around.
\end{grammarlens}

\subsection*{Your turn}
\begin{itemize}
\raggedright
  \item \emph{Say it:} \emph{kāichuán}
  \item \emph{Say it:} \emph{kāichuán}, once more
  \item \emph{Say it:} \emph{chuánzhǎng shuō shí diǎn kāichuán} — The captain says we sail at ten.
  \item \emph{Recall:} say the Mandarin for a pier for boats, then the Mandarin for a ship's captain, then say \emph{kāichuán} again
\end{itemize}

\begin{tcolorbox}[breakable,colback=teal!4,colframe=teal!35!black,title={Before you move on}]
What does \zh{开船} mean? (\textquotedblleft{}To set sail\textquotedblright{}.) Say it once more.
\end{tcolorbox}
\section[dùchuán]{dùchuán — ferry}

\label{lesson:ZH-C411-duchuan}



\begin{quote}
Before the new one: say the Mandarin for a ship's captain, then the Mandarin for to set sail.
\end{quote}

\subsection*{You'll want to know: dùchuán}
\emph{dùchuán} — \textquotedblleft{}ferry\textquotedblright{}.

\begin{quote}
\emph{wǒmen zuò dùchuán qù nàr} — We're taking the ferry over there.
\end{quote}

Its characters come later; for now it is a word for the ear.

\begin{cousinweb}
\emph{dùchuán} is built from \emph{dù} 'to cross water' and \zh{船} 'a boat', both inherited from Middle Chinese, \emph{duH} and \emph{zywen} (Guangyun, in Baxter's transcription), which glosses the first 'to cross over'.
\end{cousinweb}

\begin{grammarlens}[title={What you've built}]
One more word for getting around.
\end{grammarlens}

\subsection*{Your turn}
\begin{itemize}
\raggedright
  \item \emph{Say it:} \emph{dùchuán}
  \item \emph{Say it:} \emph{dùchuán}, once more
  \item \emph{Say it:} \emph{wǒmen zuò dùchuán qù nàr} — We're taking the ferry over there.
  \item \emph{Recall:} say the Mandarin for a ship's captain, then the Mandarin for to set sail, then say \emph{dùchuán} again
\end{itemize}

\begin{tcolorbox}[breakable,colback=teal!4,colframe=teal!35!black,title={Before you move on}]
What does dùchuán mean? (\textquotedblleft{}Ferry\textquotedblright{}.) Say it once more.
\end{tcolorbox}
\section[\texorpdfstring{\zh{救生衣} (jiùshēngyī)}{救生衣 (jiùshēngyī)}]{\zh{救生衣} (jiùshēngyī) — life jacket}

\label{lesson:ZH-C411-jiushengyi}



\begin{quote}
Before the new one: say the Mandarin for to set sail, then the Mandarin for a ferry across the water.
\end{quote}

\subsection*{You'll want to know: \zh{救生衣}}
\textbf{\zh{救生衣}} — \emph{jiùshēngyī} — \textquotedblleft{}life jacket\textquotedblright{}.

\begin{quote}
\emph{zuò chuán de shíhou yào chuān jiùshēngyī} — On a boat you have to wear a life jacket.
\end{quote}

\begin{cousinweb}
\textbf{\zh{救生衣}} (\emph{jiùshēngyī}) is built from \emph{jiù} 'to save', \zh{生} 'life' and \zh{衣} 'clothes', all inherited from Middle Chinese, \emph{kjuwH}, \emph{srjaeng} and \emph{'j+j} (Guangyun, in Baxter's transcription), which glosses the first 'to protect'.
\end{cousinweb}

\begin{grammarlens}[title={What you've built}]
One more word for getting around.
\end{grammarlens}

\subsection*{Your turn}
\begin{itemize}
\raggedright
  \item \emph{Say it:} \emph{jiùshēngyī}
  \item \emph{Say it:} \emph{jiùshēngyī}, once more
  \item \emph{Say it:} \emph{zuò chuán de shíhou yào chuān jiùshēngyī} — On a boat you have to wear a life jacket.
  \item \emph{Recall:} say the Mandarin for to set sail, then the Mandarin for a ferry across the water, then say \emph{jiùshēngyī} again
\end{itemize}

\begin{tcolorbox}[breakable,colback=teal!4,colframe=teal!35!black,title={Before you move on}]
What does \zh{救生衣} mean? (\textquotedblleft{}Life jacket\textquotedblright{}.) Say it once more.
\end{tcolorbox}
